\begin{solapartat}
A partir de la gràfica:
\figura[0.4]{sol_fig1.png}
es pot concloure que:
\begin{enumerate}[label=\arabic*-]
\item L'amplitud és: $A = 12\ \mathrm{cm}$ \punts{0,2}
\item El període és: $T = 2\times 3 = 6,0\ \mathrm{s}$ \punts{0,2}\\
i la freqüència angular és: $\omega = \frac{2\pi}{T} = \frac{\pi}{3} = 1,0\ \mathrm{rad/s}$ \punts{0,2}
\item La fase inicial és:
\[ x(t) = A\sin(\omega t + \phi_0) \Rightarrow 0 = 12\sin(\phi_0) \Rightarrow \sin(\phi_0) = 0,0 \Rightarrow \phi_0 = 0,0 \ \text{\punts{0,4}} \]
% DUBTE: amb la funció cosinus, x(t) = A cos(ωt + φ0), la fase inicial hauria de ser -π/2 (no π/2). Es transcriu tal com és.
(en el cas que facin servir la funció cosinus, la fase inicial ha de ser: $\frac{\pi}{2}$)
\end{enumerate}
\end{solapartat}

\begin{solapartat}
L'equació del moviment serà:
\[ x(t) = 12\sin\left(\frac{\pi}{3}t\right)\ \mathrm{cm} \ \text{\punts{0,4}} \left(\text{ó } x(t) = 12\cos\left(\frac{\pi}{3}t + \frac{\pi}{2}\right)\right) \]
% DUBTE: 12 cos(πt/3 + π/2) = -12 sin(πt/3); l'expressió equivalent correcta seria 12 cos(πt/3 - π/2). Es transcriu tal com és.
La constant de la molla ve donada per l'expressió:
\[ K = m\,\omega^2 = 0,25\times\left(\frac{\pi}{3}\right)^2 = 2,7\cdot 10^{-1}\ \mathrm{N/m} \ \text{\punts{0,3}} \]
L'energia mecànica del cos és:
\[ E = \frac{1}{2}\,K\,A^2 = 1,9\times 10^{-3}\ \mathrm{J} \ \text{\punts{0,3}} \]
\end{solapartat}
